%=============================================================================
% EXPERIMENT 2 APPENDIX SECTION -- copy and paste this whole block.
%
% Self-contained: no custom macros, no \input, nothing to define elsewhere.
%
% Requires in your preamble:
%   \usepackage{graphicx}
%   \usepackage{booktabs}
%   \usepackage{placeins}   % for \FloatBarrier
%
% Figure paths read "figures/experiment-2/<family>/<plot>_<condition>.pdf",
% relative to the main document. Copy the family folders from
% outputs/finals/experiment-2/{ddm,rdm,cdm}/ into <maindoc>/figures/experiment-2/
% (or point \graphicspath there). To relocate, find-and-replace the single
% prefix "figures/experiment-2/" throughout this block.
%
% Source: frozen experiment-2 snapshot, assembled 2026-09-03 by
% experiments/finalize_experiment_2.py. See README.md and MANIFEST.csv.
%=============================================================================

% --- Float placement. Defaults restored at the end of this block, so pasting
% --- it mid-document does not affect other sections.
\renewcommand{\topfraction}{0.92}
\renewcommand{\bottomfraction}{0.7}
\renewcommand{\textfraction}{0.06}
\renewcommand{\floatpagefraction}{0.6}

%=============================================================================
\section{Experiment 2: posterior predictive checks on real data}
\label{app:experiment-2-figures}
%=============================================================================

This appendix collects the posterior predictive checks for the three trained
families applied to real datasets. Posteriors were obtained from the
model-class CogFormer and used to simulate predictive datasets, which are then
compared against the observed data along two dimensions: reaction time, and the
response itself.

For every family we report three reaction-time diagnostics --- the predictive
density (\texttt{rt\_ppc}), the quantile comparison (\texttt{rt\_quantiles}),
and the quantile--quantile plot (\texttt{rt\_qq}) --- under each of two
grouping variables, listed in Table~\ref{tab:e2-conditions}. That gives
$3 \times 2 \times 3 = 18$ reaction-time panels, in
Sections~\ref{sec:e2-ddm}--\ref{sec:e2-cdm}.

Section~\ref{sec:e2-cdm-choice} then reports the response-dimension checks for
CDM: the predictive density of the response angle and its quantiles, under the
same two grouping variables, for four further panels. These are shown for CDM
only, and deliberately so. CDM's response is a continuous angle, so a density
and quantiles of $\theta$ are meaningful in the same way reaction-time
quantiles are. DDM and RDM respond with a binary choice, for which quantiles
carry no information; their response-dimension check is a comparison of choice
proportions, reported in the main text rather than here.

All predictive bands are 95\% intervals, and reaction-time quantiles are taken
at $0.1, 0.3, 0.5, 0.7, 0.9$ throughout.

\begin{table}[htbp]
  \centering
  \small
  \begin{tabular}{@{}lll@{}}
    \toprule
    Family & Grouping variable & Levels \\
    \midrule
    DDM & Stimulus type   & Novel, Familiar \\
        & Factual truth   & False, True \\
    \addlinespace
    RDM & Congruency      & Incongruent, Congruent \\
        & Session         & Session 0, Session 1 \\
    \addlinespace
    CDM & Speed instruction & Unspeeded, Speeded \\
        & Difficulty        & Easy, Medium, Hard \\
    \bottomrule
  \end{tabular}
  \caption{Grouping variables used to condition the posterior predictive
    checks, by family. Each panel of a figure corresponds to one level.}
  \label{tab:e2-conditions}
\end{table}

Two caveats travel with these figures. First, the RDM panels were produced
before the Stroop dataset landed and will need regenerating once that analysis
is settled. Second, the figures were drawn from a restyle of the plotting code
that was uncommitted at the time the snapshot was frozen: per-family identity
colours, a single 95\% interval level in place of a mix of 90\% and 95\% bands,
and display names for the grouping variables. A clean checkout reproduces the
same numbers --- the posterior draws are stored alongside the figures and are
unchanged --- but not the same appearance.

\FloatBarrier
\clearpage

%-----------------------------------------------------------------------------
\subsection{Drift-diffusion model (DDM): reaction-time checks}
\label{sec:e2-ddm}
%-----------------------------------------------------------------------------

\begin{figure}[htbp]
  \centering
  \includegraphics[width=\linewidth]{figures/experiment-2/ddm/rt_ppc_stim_type.pdf}
  \caption[DDM, stimulus type: RT posterior predictive]{\textbf{DDM, by stimulus type --- RT posterior predictive.} Observed reaction-time density (kernel density estimate) against the
    posterior predictive density, one panel per level of stimulus type (Novel and Familiar).
    The shaded band is the 95\% posterior predictive interval; thin lines are
    individual predictive draws. Misfit shows as the observed curve leaving the
    band.}
  \label{fig:e2-ddm-rt_ppc-stim_type}
\end{figure}

\begin{figure}[htbp]
  \centering
  \includegraphics[width=\linewidth]{figures/experiment-2/ddm/rt_quantiles_stim_type.pdf}
  \caption[DDM, stimulus type: RT quantiles]{\textbf{DDM, by stimulus type --- RT quantiles.} Observed reaction-time quantiles ($0.1, 0.3, 0.5, 0.7, 0.9$) by level of
    stimulus type (Novel and Familiar), each against its own 95\% posterior predictive
    interval. Complements the density check: a density can look acceptable while
    specific quantiles are systematically displaced, because the eye integrates
    over the curve.}
  \label{fig:e2-ddm-rt_quantiles-stim_type}
\end{figure}

\begin{figure}[htbp]
  \centering
  \includegraphics[width=\linewidth]{figures/experiment-2/ddm/rt_qq_stim_type.pdf}
  \caption[DDM, stimulus type: RT quantile--quantile]{\textbf{DDM, by stimulus type --- RT quantile--quantile.} Observed reaction-time quantile against the predicted quantile with a
    $45^\circ$ reference line, by level of stimulus type (Novel and Familiar). Deviation from
    the diagonal is the misfit in seconds, so systematic over- or
    under-prediction reads as a tilt rather than as two curves to subtract
    mentally.}
  \label{fig:e2-ddm-rt_qq-stim_type}
\end{figure}

\begin{figure}[htbp]
  \centering
  \includegraphics[width=\linewidth]{figures/experiment-2/ddm/rt_ppc_truth.pdf}
  \caption[DDM, factual truth: RT posterior predictive]{\textbf{DDM, by factual truth --- RT posterior predictive.} Observed reaction-time density (kernel density estimate) against the
    posterior predictive density, one panel per level of factual truth (False and True).
    The shaded band is the 95\% posterior predictive interval; thin lines are
    individual predictive draws. Misfit shows as the observed curve leaving the
    band.}
  \label{fig:e2-ddm-rt_ppc-truth}
\end{figure}

\begin{figure}[htbp]
  \centering
  \includegraphics[width=\linewidth]{figures/experiment-2/ddm/rt_quantiles_truth.pdf}
  \caption[DDM, factual truth: RT quantiles]{\textbf{DDM, by factual truth --- RT quantiles.} Observed reaction-time quantiles ($0.1, 0.3, 0.5, 0.7, 0.9$) by level of
    factual truth (False and True), each against its own 95\% posterior predictive
    interval. Complements the density check: a density can look acceptable while
    specific quantiles are systematically displaced, because the eye integrates
    over the curve.}
  \label{fig:e2-ddm-rt_quantiles-truth}
\end{figure}

\begin{figure}[htbp]
  \centering
  \includegraphics[width=\linewidth]{figures/experiment-2/ddm/rt_qq_truth.pdf}
  \caption[DDM, factual truth: RT quantile--quantile]{\textbf{DDM, by factual truth --- RT quantile--quantile.} Observed reaction-time quantile against the predicted quantile with a
    $45^\circ$ reference line, by level of factual truth (False and True). Deviation from
    the diagonal is the misfit in seconds, so systematic over- or
    under-prediction reads as a tilt rather than as two curves to subtract
    mentally.}
  \label{fig:e2-ddm-rt_qq-truth}
\end{figure}

\FloatBarrier
\clearpage

%-----------------------------------------------------------------------------
\subsection{Racing diffusion model (RDM): reaction-time checks}
\label{sec:e2-rdm}
%-----------------------------------------------------------------------------

\begin{figure}[htbp]
  \centering
  \includegraphics[width=\linewidth]{figures/experiment-2/rdm/rt_ppc_congruency.pdf}
  \caption[RDM, congruency: RT posterior predictive]{\textbf{RDM, by congruency --- RT posterior predictive.} Observed reaction-time density (kernel density estimate) against the
    posterior predictive density, one panel per level of congruency (Incongruent and Congruent).
    The shaded band is the 95\% posterior predictive interval; thin lines are
    individual predictive draws. Misfit shows as the observed curve leaving the
    band.}
  \label{fig:e2-rdm-rt_ppc-congruency}
\end{figure}

\begin{figure}[htbp]
  \centering
  \includegraphics[width=\linewidth]{figures/experiment-2/rdm/rt_quantiles_congruency.pdf}
  \caption[RDM, congruency: RT quantiles]{\textbf{RDM, by congruency --- RT quantiles.} Observed reaction-time quantiles ($0.1, 0.3, 0.5, 0.7, 0.9$) by level of
    congruency (Incongruent and Congruent), each against its own 95\% posterior predictive
    interval. Complements the density check: a density can look acceptable while
    specific quantiles are systematically displaced, because the eye integrates
    over the curve.}
  \label{fig:e2-rdm-rt_quantiles-congruency}
\end{figure}

\begin{figure}[htbp]
  \centering
  \includegraphics[width=\linewidth]{figures/experiment-2/rdm/rt_qq_congruency.pdf}
  \caption[RDM, congruency: RT quantile--quantile]{\textbf{RDM, by congruency --- RT quantile--quantile.} Observed reaction-time quantile against the predicted quantile with a
    $45^\circ$ reference line, by level of congruency (Incongruent and Congruent). Deviation from
    the diagonal is the misfit in seconds, so systematic over- or
    under-prediction reads as a tilt rather than as two curves to subtract
    mentally.}
  \label{fig:e2-rdm-rt_qq-congruency}
\end{figure}

\begin{figure}[htbp]
  \centering
  \includegraphics[width=\linewidth]{figures/experiment-2/rdm/rt_ppc_session.pdf}
  \caption[RDM, session: RT posterior predictive]{\textbf{RDM, by session --- RT posterior predictive.} Observed reaction-time density (kernel density estimate) against the
    posterior predictive density, one panel per level of session (Session 0 and Session 1).
    The shaded band is the 95\% posterior predictive interval; thin lines are
    individual predictive draws. Misfit shows as the observed curve leaving the
    band.}
  \label{fig:e2-rdm-rt_ppc-session}
\end{figure}

\begin{figure}[htbp]
  \centering
  \includegraphics[width=\linewidth]{figures/experiment-2/rdm/rt_quantiles_session.pdf}
  \caption[RDM, session: RT quantiles]{\textbf{RDM, by session --- RT quantiles.} Observed reaction-time quantiles ($0.1, 0.3, 0.5, 0.7, 0.9$) by level of
    session (Session 0 and Session 1), each against its own 95\% posterior predictive
    interval. Complements the density check: a density can look acceptable while
    specific quantiles are systematically displaced, because the eye integrates
    over the curve.}
  \label{fig:e2-rdm-rt_quantiles-session}
\end{figure}

\begin{figure}[htbp]
  \centering
  \includegraphics[width=\linewidth]{figures/experiment-2/rdm/rt_qq_session.pdf}
  \caption[RDM, session: RT quantile--quantile]{\textbf{RDM, by session --- RT quantile--quantile.} Observed reaction-time quantile against the predicted quantile with a
    $45^\circ$ reference line, by level of session (Session 0 and Session 1). Deviation from
    the diagonal is the misfit in seconds, so systematic over- or
    under-prediction reads as a tilt rather than as two curves to subtract
    mentally.}
  \label{fig:e2-rdm-rt_qq-session}
\end{figure}

\FloatBarrier
\clearpage

%-----------------------------------------------------------------------------
\subsection{Circular diffusion model (CDM): reaction-time checks}
\label{sec:e2-cdm}
%-----------------------------------------------------------------------------

\begin{figure}[htbp]
  \centering
  \includegraphics[width=\linewidth]{figures/experiment-2/cdm/rt_ppc_speed.pdf}
  \caption[CDM, speed instruction: RT posterior predictive]{\textbf{CDM, by speed instruction --- RT posterior predictive.} Observed reaction-time density (kernel density estimate) against the
    posterior predictive density, one panel per level of speed instruction (Unspeeded and Speeded).
    The shaded band is the 95\% posterior predictive interval; thin lines are
    individual predictive draws. Misfit shows as the observed curve leaving the
    band.}
  \label{fig:e2-cdm-rt_ppc-speed}
\end{figure}

\begin{figure}[htbp]
  \centering
  \includegraphics[width=\linewidth]{figures/experiment-2/cdm/rt_quantiles_speed.pdf}
  \caption[CDM, speed instruction: RT quantiles]{\textbf{CDM, by speed instruction --- RT quantiles.} Observed reaction-time quantiles ($0.1, 0.3, 0.5, 0.7, 0.9$) by level of
    speed instruction (Unspeeded and Speeded), each against its own 95\% posterior predictive
    interval. Complements the density check: a density can look acceptable while
    specific quantiles are systematically displaced, because the eye integrates
    over the curve.}
  \label{fig:e2-cdm-rt_quantiles-speed}
\end{figure}

\begin{figure}[htbp]
  \centering
  \includegraphics[width=\linewidth]{figures/experiment-2/cdm/rt_qq_speed.pdf}
  \caption[CDM, speed instruction: RT quantile--quantile]{\textbf{CDM, by speed instruction --- RT quantile--quantile.} Observed reaction-time quantile against the predicted quantile with a
    $45^\circ$ reference line, by level of speed instruction (Unspeeded and Speeded). Deviation from
    the diagonal is the misfit in seconds, so systematic over- or
    under-prediction reads as a tilt rather than as two curves to subtract
    mentally.}
  \label{fig:e2-cdm-rt_qq-speed}
\end{figure}

\begin{figure}[htbp]
  \centering
  \includegraphics[width=\linewidth]{figures/experiment-2/cdm/rt_ppc_difficulty.pdf}
  \caption[CDM, difficulty: RT posterior predictive]{\textbf{CDM, by difficulty --- RT posterior predictive.} Observed reaction-time density (kernel density estimate) against the
    posterior predictive density, one panel per level of difficulty (Easy, Medium and Hard).
    The shaded band is the 95\% posterior predictive interval; thin lines are
    individual predictive draws. Misfit shows as the observed curve leaving the
    band.}
  \label{fig:e2-cdm-rt_ppc-difficulty}
\end{figure}

\begin{figure}[htbp]
  \centering
  \includegraphics[width=\linewidth]{figures/experiment-2/cdm/rt_quantiles_difficulty.pdf}
  \caption[CDM, difficulty: RT quantiles]{\textbf{CDM, by difficulty --- RT quantiles.} Observed reaction-time quantiles ($0.1, 0.3, 0.5, 0.7, 0.9$) by level of
    difficulty (Easy, Medium and Hard), each against its own 95\% posterior predictive
    interval. Complements the density check: a density can look acceptable while
    specific quantiles are systematically displaced, because the eye integrates
    over the curve.}
  \label{fig:e2-cdm-rt_quantiles-difficulty}
\end{figure}

\begin{figure}[htbp]
  \centering
  \includegraphics[width=\linewidth]{figures/experiment-2/cdm/rt_qq_difficulty.pdf}
  \caption[CDM, difficulty: RT quantile--quantile]{\textbf{CDM, by difficulty --- RT quantile--quantile.} Observed reaction-time quantile against the predicted quantile with a
    $45^\circ$ reference line, by level of difficulty (Easy, Medium and Hard). Deviation from
    the diagonal is the misfit in seconds, so systematic over- or
    under-prediction reads as a tilt rather than as two curves to subtract
    mentally.}
  \label{fig:e2-cdm-rt_qq-difficulty}
\end{figure}

\FloatBarrier
\clearpage

%-----------------------------------------------------------------------------
\subsection{Circular diffusion model (CDM): response-dimension checks}
\label{sec:e2-cdm-choice}
%-----------------------------------------------------------------------------

\begin{figure}[htbp]
  \centering
  \includegraphics[width=\linewidth]{figures/experiment-2/cdm/choice_ppc_speed.pdf}
  \caption[CDM, speed instruction: response-dimension posterior predictive]{\textbf{CDM, by speed instruction --- response-dimension posterior predictive.} Observed density of the response angle $\theta$ against the posterior
    predictive density, one panel per level of speed instruction (Unspeeded and Speeded), with the
    95\% posterior predictive interval shaded. CDM's response is a continuous
    angle, so its misfit is invisible to an RT-only check.}
  \label{fig:e2-cdm-choice_ppc-speed}
\end{figure}

\begin{figure}[htbp]
  \centering
  \includegraphics[width=\linewidth]{figures/experiment-2/cdm/choice_quantiles_speed.pdf}
  \caption[CDM, speed instruction: response-dimension quantiles]{\textbf{CDM, by speed instruction --- response-dimension quantiles.} Quantiles ($0.1, 0.3, 0.5, 0.7, 0.9$) of the observed response angle
    $\theta$ against their posterior predictive intervals, by level of speed instruction
    (Unspeeded and Speeded). Meaningful for CDM because the response is continuous; the
    binary-choice families have no counterpart.}
  \label{fig:e2-cdm-choice_quantiles-speed}
\end{figure}

\begin{figure}[htbp]
  \centering
  \includegraphics[width=\linewidth]{figures/experiment-2/cdm/choice_ppc_difficulty.pdf}
  \caption[CDM, difficulty: response-dimension posterior predictive]{\textbf{CDM, by difficulty --- response-dimension posterior predictive.} Observed density of the response angle $\theta$ against the posterior
    predictive density, one panel per level of difficulty (Easy, Medium and Hard), with the
    95\% posterior predictive interval shaded. CDM's response is a continuous
    angle, so its misfit is invisible to an RT-only check.}
  \label{fig:e2-cdm-choice_ppc-difficulty}
\end{figure}

\begin{figure}[htbp]
  \centering
  \includegraphics[width=\linewidth]{figures/experiment-2/cdm/choice_quantiles_difficulty.pdf}
  \caption[CDM, difficulty: response-dimension quantiles]{\textbf{CDM, by difficulty --- response-dimension quantiles.} Quantiles ($0.1, 0.3, 0.5, 0.7, 0.9$) of the observed response angle
    $\theta$ against their posterior predictive intervals, by level of difficulty
    (Easy, Medium and Hard). Meaningful for CDM because the response is continuous; the
    binary-choice families have no counterpart.}
  \label{fig:e2-cdm-choice_quantiles-difficulty}
\end{figure}

\FloatBarrier

% --- Restore LaTeX default float parameters -----------------------------------
\renewcommand{\topfraction}{0.7}
\renewcommand{\bottomfraction}{0.3}
\renewcommand{\textfraction}{0.2}
\renewcommand{\floatpagefraction}{0.5}

%=============================================================================
% END OF EXPERIMENT 2 APPENDIX SECTION
%=============================================================================
